%
%Academic CV LaTeX Template
% Author: Dubasi Pavan Kumar
% LinkedIn: https://www.linkedin.com/in/im-pavankumar/
% License: MIT
%
% For errors, suggestions, or improvements, please contact:
% Email: pavankumard.pg19.ma@nitp.ac.in
%



\documentclass[a4paper,11pt]{article}

% Configuración para mejor copia de texto
\usepackage{ifxetex,ifluatex}
\ifxetex
  \usepackage{fontspec}
  \setmainfont{Latin Modern Roman}
\else
  \ifluatex
    \usepackage{fontspec}
    \setmainfont{Latin Modern Roman}
  \else
    \usepackage{cmap}
    \usepackage[T1]{fontenc}
    \usepackage{lmodern}
    \input glyphtounicode.tex
    \pdfgentounicode=1
  \fi
\fi


% Package imports
\usepackage{latexsym}
\usepackage{xcolor}
\usepackage{float}
\usepackage{ragged2e}
\usepackage[empty]{fullpage}
\usepackage{wrapfig}
\usepackage{lipsum}
\usepackage{tabularx}
\usepackage{titlesec}
\usepackage{geometry}
\usepackage{marvosym}
\usepackage{verbatim}
\usepackage{enumitem}
\usepackage{fancyhdr}
\usepackage{multicol}
\usepackage{graphicx}
\usepackage{cfr-lm}

\usepackage{fontawesome5}

\usepackage{microtype}
\usepackage[utf8]{inputenc}
\usepackage[spanish]{babel}

% Configuración mejorada para copia de texto
\usepackage{selnolig}
\usepackage{accsupp}

% Configuración específica para mejorar copia de texto

% Configuración de microtype más agresiva
\microtypesetup{
    tracking=true,
    kerning=true,
    spacing=true,
    factor=1100,
    stretch=20,
    shrink=20
}

% Evitar saltos de línea problemáticos
\tolerance=1414
\hbadness=1414
\emergencystretch=1.5em
\hfuzz=0.3pt
\widowpenalty=9999
\vfuzz=\hfuzz
\raggedbottom



% Color definitions
\definecolor{darkblue}{RGB}{0,0,139}

% Page layout
\setlength{\multicolsep}{0pt} 
\pagestyle{fancy}
\fancyhf{} % clear all header and footer fields
\fancyfoot{}
\renewcommand{\headrulewidth}{0pt}
\renewcommand{\footrulewidth}{0pt}
\geometry{left=1.4cm, top=0.8cm, right=1.2cm, bottom=1cm}
\setlength{\footskip}{5pt} % Addressing fancyhdr warning

% Hyperlink setup (moved after fancyhdr to address warning)
\usepackage[hidelinks]{hyperref}
\hypersetup{
    colorlinks=true,
    linkcolor=darkblue,
    filecolor=darkblue,
    urlcolor=darkblue,
    pdfencoding=auto,
    psdextra,
    pdfpagemode=UseNone,
    pdfstartview=FitH,
    pdfcreator={LaTeX},
    pdfproducer={pdfTeX},
    pdftex,
    unicode=true
}


% Custom box settings
\usepackage[most]{tcolorbox}
\tcbset{
    frame code={},
    center title,
    left=0pt,
    right=0pt,
    top=0pt,
    bottom=0pt,
    colback=gray!20,
    colframe=white,
    width=\dimexpr\textwidth\relax,
    enlarge left by=-2mm,
    boxsep=4pt,
    arc=0pt,outer arc=0pt,
}

% URL style
\urlstyle{same}

% Text alignment
\setlength{\tabcolsep}{0in}

% Section formatting
\titleformat{\section}{
  \vspace{-4pt}\scshape\large
}{}{0em}{}[\color{black}\titlerule \vspace{-7pt}]

% Custom commands
\newcommand{\resumeItem}[2]{
  \item{
    \textbf{#1}{\hspace{0.5mm}#2 \vspace{-0.5mm}}
  }
}

\newcommand{\resumePOR}[3]{
\vspace{0.5mm}\item
    \begin{tabular*}{0.97\textwidth}[t]{l@{\extracolsep{\fill}}r}
        \textbf{#1}\hspace{0.3mm}#2 & \textit{\small{#3}} 
    \end{tabular*}
    \vspace{-2mm}
}

\newcommand{\resumeSubheading}[4]{
\vspace{0.5mm}\item
    \begin{tabular*}{0.98\textwidth}[t]{l@{\extracolsep{\fill}}r}
        \textbf{#1} & \textit{\footnotesize{#4}} \\
        \textit{\footnotesize{#3}} &  \footnotesize{#2}\\
    \end{tabular*}
    \vspace{-2.4mm}
}

\newcommand{\resumeProject}[4]{
\vspace{0.5mm}\item
    \begin{tabular*}{0.98\textwidth}[t]{l@{\extracolsep{\fill}}r}
        \textbf{#1} & \textit{\footnotesize{#3}} \\
        \footnotesize{\textit{#2}} & \footnotesize{#4}
    \end{tabular*}
    \vspace{-2.4mm}
}

\newcommand{\resumeSubItem}[2]{\resumeItem{#1}{#2}\vspace{-4pt}\footnotesize}

\renewcommand{\labelitemi}{$\vcenter{\hbox{\tiny$\bullet$}}$}
\renewcommand{\labelitemii}{$\vcenter{\hbox{\tiny$\circ$}}$}

\newcommand{\resumeSubHeadingListStart}{\begin{itemize}[leftmargin=*,labelsep=1mm]}
\newcommand{\resumeHeadingSkillStart}{\begin{itemize}[leftmargin=*,itemsep=1.7mm, rightmargin=2ex]}
\newcommand{\resumeItemListStart}{\begin{itemize}[leftmargin=*,labelsep=1mm,itemsep=0.5mm]\footnotesize}

\newcommand{\resumeSubHeadingListEnd}{\end{itemize}\vspace{2mm}}
\newcommand{\resumeHeadingSkillEnd}{\end{itemize}\vspace{-2mm}}
\newcommand{\resumeItemListEnd}{\end{itemize}\vspace{-2mm}}
\newcommand{\cvsection}[1]{%
\vspace{2mm}
\begin{tcolorbox}
    \textbf{\large #1}
\end{tcolorbox}
    \vspace{-4mm}
}

\newcolumntype{L}{>{\arraybackslash}X}%
\newcolumntype{R}{>{\raggedleft\arraybackslash}X}%
\newcolumntype{C}{>{\centering\arraybackslash}X}%

% Commands for icon sizing and positioning
\newcommand{\socialicon}[1]{\raisebox{-0.05em}{\resizebox{!}{1em}{#1}}}
\newcommand{\ieeeicon}[1]{\raisebox{-0.3em}{\resizebox{!}{1.3em}{#1}}}

% Macro para copiar párrafos sin saltos de línea añadidos por el visor PDF
\newcommand{\copyablepar}[1]{%
\BeginAccSupp{method=escape,ActualText={#1}}#1\EndAccSupp{}
}

% Font options
\newcommand{\headerfonti}{\fontfamily{phv}\selectfont} % Helvetica-like (similar to Arial/Calibri)
\newcommand{\headerfontii}{\fontfamily{ptm}\selectfont} % Times-like (similar to Times New Roman)
\newcommand{\headerfontiii}{\fontfamily{ppl}\selectfont} % Palatino (elegant serif)
\newcommand{\headerfontiv}{\fontfamily{pbk}\selectfont} % Bookman (readable serif)
\newcommand{\headerfontv}{\fontfamily{pag}\selectfont} % Avant Garde-like (similar to Trebuchet MS)
\newcommand{\headerfontvi}{\fontfamily{cmss}\selectfont} % Computer Modern Sans Serif
\newcommand{\headerfontvii}{\fontfamily{qhv}\selectfont} % Quasi-Helvetica (another Arial/Calibri alternative)
\newcommand{\headerfontviii}{\fontfamily{qpl}\selectfont} % Quasi-Palatino (another elegant serif option)
\newcommand{\headerfontix}{\fontfamily{qtm}\selectfont} % Quasi-Times (another Times New Roman alternative)
\newcommand{\headerfontx}{\fontfamily{bch}\selectfont} % Charter (clean serif font)

\begin{document}
\headerfontiii

% Configuración adicional para mejorar copia de texto

% Header
\begin{center}
    {\huge\textbf{Castro Elvira Diego}}
\end{center}
\vspace{-6mm}

\begin{center}
    \small{
   Ingeniero en Inteligencia Artificial y Científico de Datos
    }
\end{center}
\vspace{-6mm}

\begin{center}
    \small{
    +52 2411081478 | \href{diego.castro.elvira@gmail.com}{diego.castro.elvira@gmail.com} | 
    \socialicon{\faLinkedin} \href{https://www.linkedin.com/in/diego-castro-elvira/}{diego-castro-elvira} | 
    \socialicon{\faGithub} \href{https://github.com/DiegoCastr00}{DiegoCastr00} 
    }
\end{center}


\vspace{-4mm}

\section{\textbf{Perfil Profesional}}
\vspace{1mm}
\footnotesize{
Ingeniero en Inteligencia Artificial especializado en ciencia de datos, con experiencia en visión computacional, NLP y modelado de series de tiempo. Experto en el ciclo de vida completo de modelos: desde el procesamiento de datos hasta el entrenamiento, evaluación y despliegue en la nube. Aplicación de técnicas avanzadas de deep learning (fine-tuning, aprendizaje por transferencia, aprendizaje contrastivo) y modelado estadístico para resolver problemas complejos con impacto real.}
\vspace{-2mm}



\section{\textbf{Experiencia Profesional}}
\vspace{-0.4mm}
  \resumeSubHeadingListStart
  \resumeSubheading
      {{LatGotLab [\href{https://latgoblab.com/}{\faIcon{globe}}]}}{México, híbrido}
      {{ Full-Stack AI Engineer }}{Julio 2024 – Agosto 2025}
      \resumeItemListStart
        \item Desarrollé una biblioteca de Recuperación Aumentada con Generación (RAG) para consulta inteligente de libros y documentos, integrando modelos de lenguaje para responder en lenguaje natural.
        \item Lidero el desarrollo de Noosfera [\href{https://noosfera.ai/}{\faIcon{globe}}], aplicación móvil en React Native con clasificación automática de imágenes mediante modelos de IA, geolocalización y visualización de eventos en un mapa interactivo.
        \item Implementé un pipeline completo de IA para descripción y categorización de imágenes, conectado a una base de datos y APIs para su publicación en un feed dinámico priorizado por proximidad geográfica.
      \resumeItemListEnd 
  \resumeSubheading
    {Centro de Innovación y Desarrollo Tecnológico en Cómputo, IPN [\href{https://www.cidetec.ipn.mx/}{\faIcon{globe}}]}{Remoto}
    {Asistente de investigación}{Abril 2025 - Julio 2025}
    \resumeItemListStart
    \item Etiqueté imágenes con Roboflow para modelos de detección de objetos (\texttt{person}, \texttt{backpack}, \texttt{bag}).
      \item Aseguré la consistencia y calidad del dataset, optimizando su preparación para entrenamiento de modelos de visión por computadora.
    \resumeItemListEnd
  \resumeSubHeadingListEnd
\vspace{-6mm}

\section{\textbf{Proyectos Destacados}}
\vspace{-0.4mm}
\resumeSubHeadingListStart

\resumeProject
  {Modelo de similitud entre arte original y generado con IA}
  {Tools: PyTorch, Transformers, BLIP-2, LLaMA 3, Stable Diffusion, OpenCV, Pandas, CUDA, Multiprocessing}
  {Diciembre 2024 - Abril 2025}
  {{}\href{https://github.com/DiegoCastr00/CalculoSimilitud/tree/master}{\textcolor{darkblue}{\faGithub}}}
\resumeItemListStart
  \item Implementé una red siamesa basada en CLIP y capas proyectivas para medir similitud semántico-visual entre arte humano y generado con Stable Diffusion XL.
  \item Creé un dataset de 80k imágenes con descripciones BLIP-2 y prompts diversificados con LLaMA 3 y BART.
  \item Diseñé un pipeline de entrenamiento con estrategia de tripletas, pérdida contrastiva e InfoNCE, evaluado mediante Triplet Accuracy y métricas de similitud coseno en entorno multi-GPU, logrando alta discriminación entre pares similares y no relacionados.
\resumeItemListEnd

\resumeProject
  {Detección de emociones en imágenes}
  {Tools: Python, OpenCV, TensorFlow/Keras, NumPy, Matplotlib, CK+, FER2013.}
  {Abril 2024 – Junio 2024}
  {{}\href{https://github.com/DiegoCastr00/DeteccionEmociones}{\textcolor{darkblue}{\faGithub}}}
\resumeItemListStart
  \item Implementé un modelo de clasificación de emociones faciales (ira, disgusto, miedo, felicidad, tristeza, sorpresa, neutralidad) usando CNNs y descriptores HOG, SIFT y LBP.
  \item Entrené y evalué el modelo con datasets CK+ y FER2013, alcanzando hasta 98.47\% de precisión en entornos controlados (CK+). Analicé el impacto del balance y consistencia de datos en el rendimiento, identificando sesgos y problemas de generalización.
\resumeItemListEnd

\resumeProject
  {Análisis y Simulación de Series Temporales con Modelos ARMA}
  {Tools: Python, NumPy, Pandas, Matplotlib, Statsmodels}
  {Enero 2024 – Febrero 2024}
  {{}\href{https://github.com/DiegoCastr00/Machine_Learning/blob/master/Practica_9_ARMA/ARMA.ipynb}{\textcolor{darkblue}{\faGithub}}}
\resumeItemListStart
  \item Desarrollé y simulé series temporales utilizando modelos Autoregresivos (AR) y de Media Móvil (MA) de órdenes 1 y 2.
  \item Analicé sistemáticamente el impacto de los coeficientes del modelo ($\phi$ y $\theta$) en el comportamiento de las series, evaluando su efecto en la volatilidad, tendencia y persistencia de la autocorrelación.
  \item Utilicé funciones de autocorrelación (ACF) para interpretar la estructura de dependencia temporal, validando cómo los parámetros de los modelos AR y MA determinan los patrones subyacentes en los datos.
\resumeItemListEnd


\section{\textbf{Educación}}

\resumeProject
{Unidad Profesional Interdisciplinaria en Ingeniería Campus Tlaxcala IPN}{Tlaxcala, México}
{Ingeniería en Inteligencia Artificial}{Junio 2021 - Junio 2025}


\vspace{-3mm}

\resumeSubHeadingListEnd


\section{\textbf{Competencias y Habilidades Técnicas}}
\vspace{-0.4mm}
\resumeHeadingSkillStart

%--- IDIOMAS Y LOGROS (Ligeramente ajustado para mayor claridad) ---
\resumeSubItem{Idiomas:}
  {Español (Nativo) · Inglés (B2 - Profesional Intermedio, lectura técnica fluida)}

\resumeSubItem{Logros Clave:}
  {1° lugar – Hackathon Interno IPN (2023) · 2° lugar – Mobility Hackathon (2024) · Participación destacada – Enactus México (2025)}

%--- HABILIDADES TÉCNICAS (Reorganizado para ATS y claridad) ---
\resumeSubItem{Áreas de Especialidad:}
  {Ciencia de Datos · Visión por Computadora (CV) · Procesamiento de Lenguaje Natural (NLP) · Análisis de Series Temporales · Machine Learning Ops (MLOps)}

\resumeSubItem{Modelado y Técnicas:}
  {Deep Learning (Arquitecturas Transformer, CNNs, RNNs) · Modelado Estadístico y Predictivo · Fine-Tuning y Transfer Learning · Retrieval-Augmented Generation (RAG) · Aprendizaje por Refuerzo (RL)}

\resumeSubItem{Lenguajes y Datos:}
  {Python · SQL · TypeScript · Bash · YAML}

\resumeSubItem{Librerías y Ecosistema de IA:}
  {PyTorch · TensorFlow/Keras · Hugging Face (Transformers, Datasets) · Scikit-learn · LangChain · OpenCV · Pandas · NumPy · Polars · Statsmodels}

\resumeSubItem{Plataformas y MLOps:}
  {AWS (S3, EC2, SageMaker) · Docker · Kubernetes (K8s) · Git · CI/CD (GitHub Actions) · Wandb · MLflow · Vector DBs (FAISS, ChromaDB, Pinecone)}

\resumeSubItem{Desarrollo y Visualización:}
  {FastAPI · Node.js · React · Next.js · Power BI · Matplotlib · Seaborn}

%--- HABILIDADES ESTRATÉGICAS (Enfocadas en impacto de negocio) ---
\resumeSubItem{Habilidades Estratégicas:}
  {Liderazgo Técnico y Mentoría · Gestión de Proyectos (Agile/Scrum) · Análisis Exploratorio de Datos (EDA) · Ética en IA y Gobernanza de Datos · Resolución de Problemas Complejos}

\resumeHeadingSkillEnd

\end{document}